\documentclass[11pt]{article}

\newif\ifneuripsstyle
\IfFileExists{neurips_2024.sty}{
  \neuripsstyletrue
  \usepackage[final]{neurips_2024}
}{
  \neuripsstylefalse
  \usepackage[margin=1in]{geometry}
  \usepackage[numbers]{natbib}
  \usepackage[colorlinks=true,allcolors=blue]{hyperref}
}

\usepackage[utf8]{inputenc}
\usepackage[T1]{fontenc}
\usepackage{microtype}
\usepackage{amsmath}
\usepackage{amssymb}
\usepackage{booktabs}
\usepackage{graphicx}
\usepackage{url}

\title{Execution Intervention for Post-Hoc Debugging of LLM Agent Trajectories}

\author{Anonymous Workshop Submission}

\begin{document}

\maketitle

\begin{abstract}
Tool-using language agents often fail because one intermediate action corrupts
the downstream execution state, yet standard benchmark evaluation usually
reveals only the final reward and a raw trajectory trace. We propose execution
intervention for post-hoc debugging: a methodology that treats failed agent
trajectories as replayable programs. Given a failed run with checkpointed
environment state, the method restores an intermediate snapshot, applies a
small structured patch, replays execution, and checks whether the benchmark
outcome flips from failure to success. The smallest successful intervention is
useful both as a repair and as a causal autopsy of the root failure. We
instantiate this protocol on real tau-style environments with deterministic
state, saved benchmark trajectories, and benchmark-native evaluators. The
current system supports structured patch families, bounded search-cost
accounting, strict versus oracle continuation analysis, and explicit rerun
baselines including retry from scratch and retry from a localized snapshot. On
controlled synthetic failures, the prototype supports exact localization
metrics in addition to recovery. On imported natural failures, the current
artifact-backed evidence already includes strict non-oracle recoveries in both
retail and airline settings, including continuation-based recovery after an
identity-lookup repair and deletion-based recovery for harmful mutating
actions. These results motivate replayable execution intervention as a practical
debugging protocol for LLM agents: more causal than prompt-only retry, more
interpretable than unconstrained regeneration, and naturally aligned with
benchmark evaluation.
\end{abstract}

\section{Introduction}

Many agent benchmarks now measure multi-step interaction with tools, APIs,
databases, browsers, or code environments. When such an agent fails, the usual
options are limited: inspect the trace manually, rerun with a stronger model,
or prompt the model to reflect and retry. These approaches can improve task
success, but they do not cleanly answer a debugging question: which
intermediate decision actually caused the failure?

Our central claim is that replayable environments let us ask a stronger causal
question. If we can checkpoint and restore execution state at each step, then a
failed trajectory is not merely a transcript. It is a stateful program that
can be paused, forked, edited, and re-executed. This enables a simple but
useful search problem: find the smallest intervention that flips the final
benchmark outcome.

This workshop paper is methodology-first. We do not claim benchmark-wide repair
at this stage. Instead, we contribute a concrete protocol built around
checkpointed replay, structured patch families, bounded search cost, and
autopsy generation, together with early real-environment evidence on saved
tau-style failures in retail and airline domains.

\section{Method}

For a failed trajectory, we log per-step prompt context, chosen action, tool
result, and pre/post execution snapshots. A patch candidate specifies a target
step, a patch family, and a small structured payload. Applying a patch means:
\begin{enumerate}
  \item restore the saved pre-step snapshot;
  \item edit the action or context according to the patch payload;
  \item resume execution under a bounded continuation policy;
  \item evaluate the replayed rollout with the original benchmark evaluator.
\end{enumerate}

The current system supports structured intervention families including tool
argument edits, tool-call replacement, insertion, deletion, context edits, and
continuation replacement. Search returns the smallest successful patch under a
budget over localization, proposal, and replay. The output is not only a
success flag: it is also an autopsy artifact with the root-cause step, patch
family, patch size, total search cost, and before/after state fragments.

The workshop implementation also includes three simple rerun baselines:
\begin{itemize}
  \item retry from scratch,
  \item retry from the localized snapshot without structured patching,
  \item raw continuation from the localized snapshot while preserving the
  original failing action.
\end{itemize}

\section{Workshop Evaluation Scope}

The workshop path stays tau-native and focuses on retail plus airline. The main
headline setting is strict non-oracle replay: after patching a saved failure,
the system must continue from the replayed state without reusing oracle future
actions. Oracle continuation is retained only as an upper-bound analysis.

The codebase now supports the full workshop artifact path: building natural and
synthetic corpora, running structured patch search, running retry baselines,
exporting tables and figure CSVs, and generating reviewer-facing artifact
bundles. The current repository evidence is still moderate scale rather than a
fully frozen workshop corpus, so the manuscript should present the current
numbers as early real-environment evidence and reserve larger frozen counts for
the submission run.

\section{Related Work}

Our work sits between benchmark evaluation for interactive agents and
post-hoc repair or deliberation methods for language models.

\paragraph{Interactive agent benchmarks.}
Recent benchmarks such as tau-style environments, WebArena, WorkArena,
AgentDojo, and SWE-bench make agent execution state, tools, and outcome
evaluation central to the task definition
\citep{tau_style,web_arena,work_arena,agentdojo,swe_bench}. These benchmarks
are valuable because they expose realistic multi-step failures. Our goal is not
to propose another benchmark, but to turn benchmark trajectories into
debuggable execution objects by adding checkpointed intervention and replay.

\paragraph{Reflection and repair for language models.}
Methods such as Reflexion, Self-Refine, and Tree of Thoughts improve behavior
through critique, self-correction, or explicit search
\citep{reflexion,self_refine,tree_of_thoughts}. These methods usually operate
in prompt space: they ask the model to revise, deliberate, or retry. Our
setting is different. We intervene on replayable execution state and score each
intervention by whether it causally flips the benchmark outcome. This makes the
resulting repair more directly tied to an identifiable state change.

\paragraph{Program repair and debugging perspective.}
The closest conceptual analogy is automated debugging or program repair, where
one searches for a small change that fixes downstream behavior. Our setting is
harder in some ways because the trajectory mixes language, tool calls, and
mutable environment state. At the same time, benchmark-native replay gives us a
clean reward signal and exact state restoration, which makes causal debugging
practical.

\section{Current Contribution}

What is already true in the repository:
\begin{itemize}
  \item real tau-native replay exists for retail and airline;
  \item saved natural failures are imported only if they remain failures under
  exact replay;
  \item strict non-oracle natural recovery is no longer only synthetic, with
  artifact-backed retail continuation recovery and airline deletion recovery;
  \item workshop-specific retry baselines and artifact bundling are implemented
  in the experiment harness.
\end{itemize}

What still separates the workshop version from a full conference paper:
\begin{itemize}
  \item a larger frozen natural corpus,
  \item stronger non-oracle continuation at scale,
  \item broader baseline coverage,
  \item human evaluation of autopsy usefulness.
\end{itemize}

\section{Limitations and Next Steps}

The current system is still limited by continuation quality, patch-family
coverage, and experiment scale. Some failures are easy to localize but hard to
repair without a stronger continuation proposer. Others likely require context
editing or multi-step intervention beyond the current one-step framework. The
next milestone is to freeze the workshop corpora, run the new retry baselines
under matched budgets, and finalize the short paper with saved tables and case
studies generated directly from artifact directories.

\bibliographystyle{plainnat}
\bibliography{references}

\end{document}
